% !PRINTED 531
\weberpage{531}
Die Discriminante $\varDelta(\omega_1', \omega_2', \ldots, \omega_n')$ ist also kleiner als $\varDelta(\omega_1, \omega_2, \ldots, \omega_n)$, und dies widerspricht der Annahme, dass $\omega_1, \omega_2, \ldots, \omega_n$ eine Minimalbasis sei. Damit ist unser Satz erwiesen.

Bezeichnet man die Gesammtheit aller \gesperrt{ganzen Zahlen} des Körpers $\Omega$ mit $\mathfrak{o}$, so erhält man alle Zahlen von $\mathfrak{o}$, und jede nur einmal, wenn man in
\weq{6}{k_1\omega_1 + k_2\omega_2 + \cdots + k_n\omega_n}
die Coefficienten $k_1, k_2, \ldots, k_n$ die sämmtlichen ganzen rationalen Zahlen durchlaufen lässt.

Aus diesem Grunde wird eine Minimalbasis von $\Omega$ auch eine \gesperrt{Basis von $\mathfrak{o}$} genannt.

\gesperrt{Die Discriminante einer Basis von $\mathfrak{o}$ ist gleich der Grundzahl des Körpers $\Omega$.}

Nach der Formel (11) des vorigen Paragraphen können wir aus einer Basis von $\mathfrak{o}$ beliebig viele andere durch lineare Substitution ableiten:
\weq{7}{(\omega_1', \omega_2', \ldots, \omega_n') = C(\omega_1, \omega_2, \ldots, \omega_n),}
wenn $C$ (nach den Bezeichnungen des \S.\,37) eine lineare Substitution mit rationalen ganzzahligen Coëfficienten und der Determinante $\pm 1$ bedeutet.

Da nach der Bedeutung der Basis von $\mathfrak{o}$ alle ganzen Zahlen in $\Omega$, also auch die Elemente $\omega_1', \omega_2', \ldots, \omega_n'$ einer zweiten Basis von $\mathfrak{o}$ linear mit ganzen rationalen Coëfficienten durch $\omega_1, \omega_2, \ldots, \omega_n$ darstellbar sind, so folgt auch umgekehrt, dass man durch solche lineare Substitutionen mit der Determinante $\pm 1$ aus einer Basis von $\mathfrak{o}$ alle anderen ableiten kann.

Denn ist
\[
(\omega_1, \omega_2, \ldots, \omega_n) = C'(\omega_1', \omega_2', \ldots, \omega_n'),
\]
so muss die zusammengesetzte Substitution $CC'$ die identische sein, und das Product beider Determinanten ist also 1, also jede von ihnen $= \pm 1$.

Da unter den Zahlen von $\mathfrak{o}$ immer die Zahl 1 enthalten ist, so kann man, wenn $\omega_1, \omega_2, \ldots, \omega_n$ eine Basis von $\mathfrak{o}$ ist, die \gesperrt{ganzen rationalen Zahlen} $c_1, c_2, \ldots, c_n$ so bestimmen, dass die Relation
\weq{8}{c_1\omega_1 + c_2\omega_2 + \cdots + c_n\omega_n = 1}
befriedigt ist.
